\chapter{Criterio de trazabilidad operativa}

\section{Trazabilidad disciplinada}
\begin{definition}
Se llama trazabilidad operativa a la posibilidad de reconstruir de manera legible el recorrido disciplinado de una pieza dentro del canon intermedio, identificando su origen, su dependencia, su nivel de tratamiento y sus transformaciones relevantes.
\end{definition}

\begin{observation}
La trazabilidad operativa no equivale todavía a auditoría completa. Su función es fijar una primera condición de seguimiento suficiente para evitar opacidad estructural.
\end{observation}

\section{Condición mínima de trazabilidad}
\begin{definition}
Existe trazabilidad operativa mínima cuando pueden declararse sin contradicción inmediata:
\begin{enumerate}
  \item el punto de entrada de la pieza en el canon;
  \item la base canónica de la que hereda sentido;
  \item el nivel o niveles en los que ha sido tratada;
  \item las transformaciones, traducciones o articulaciones que la afectan;
  \item la condición bajo la cual sigue siendo reconocible como la misma pieza tratable.
\end{enumerate}
\end{definition}

\begin{observation}
No se exige todavía una granularidad exhaustiva de todos los pasos posibles. Basta con una reconstrucción suficiente para sostener seguimiento disciplinado.
\end{observation}

\section{Trazabilidad y configuración}
\begin{definition}
Una configuración operativa es trazable cuando su función, su herencia y su paso por el canon intermedio pueden reconstruirse sin pérdida inmediata de identidad tratable.
\end{definition}

\begin{observation}
Esta cláusula evita que una configuración quede reducida a una mera aparición local sin historia legible dentro del corpus.
\end{observation}

\section{Trazabilidad y familia}
\begin{definition}
Una familia operativa es trazable cuando puede declararse qué rasgos comunes conserva, qué diferencias internas admite y cómo sus miembros siguen siendo reconocibles dentro de una misma gramática disciplinada.
\end{definition}

\begin{observation}
La trazabilidad de una familia no exige uniformidad absoluta entre sus miembros. Exige continuidad suficiente de pertenencia y función.
\end{observation}

\section{Trazabilidad y traducción}
\begin{definition}
Una traducción entre niveles es trazable cuando puede declararse qué se preserva, qué se modifica y bajo qué criterio la pieza traducida conserva reconocimiento suficiente en el nivel de llegada.
\end{definition}

\begin{observation}
Esta cláusula impide que la traducción entre niveles funcione como desplazamiento opaco o como sustitución sin registro.
\end{observation}

\section{Condición mínima de fallo}
\begin{definition}
Se considera fallida una trazabilidad operativa cuando una pieza ya no puede reconstruir de manera legible su origen, su dependencia, su nivel o sus transformaciones relevantes dentro del canon intermedio.
\end{definition}

\begin{observation}
Esta condición protege al Tomo II frente a piezas cuyo uso se vuelve irrepetible u oscuro por pérdida de continuidad interna.
\end{observation}

\section{Cláusula de alcance}
\begin{observation}
El criterio de trazabilidad operativa fija una primera gramática de seguimiento disciplinado dentro del canon intermedio. No constituye todavía una teoría completa de auditoría final ni de validación exhaustiva.
\end{observation}

